\documentclass[10pt]{article}
\usepackage[margin=0.8in]{geometry}
\usepackage[T1]{fontenc}
\usepackage{lmodern}
\usepackage{hyperref}
\usepackage{booktabs}
\hypersetup{hidelinks}
\setlength{\emergencystretch}{2em}
\newcommand{\Verdict}{not explained}
\newcommand{\Rcount}{10}
\newcommand{\Scount}{10}
\newcommand{\Bcount}{0}
\newcommand{\Ccount}{0}
\newcommand{\Pvalue}{1.0000}
\newcommand{\Pairs}{29}
\newcommand{\Dropped}{1}
\newcommand{\FixedSilentWitness}{0}
\newcommand{\SilentTotal}{10}
\newcommand{\FixedRows}{\texttt{dense} & 1 & 0 & 0 & 28 & 28 \\
\texttt{dense} & 2 & 0 & not launched & 27 & -- \\
\texttt{dense} & 3 & 3 & not launched & 21 & -- \\
\texttt{dense-latency-100us} & 1 & 0 & not launched & 28 & -- \\
\texttt{dense-latency-100us} & 2 & 0 & not launched & 28 & -- \\
\texttt{dense-latency-100us} & 3 & 3 & not launched & 21 & -- \\
\texttt{dense-latency-10ms} & 1 & 0 & not launched & 28 & -- \\
\texttt{dense-latency-10ms} & 2 & 7 & not launched & 13 & -- \\
\texttt{dense-latency-10ms} & 3 & 0 & not launched & 21 & -- \\
\texttt{dense-service-100us} & 1 & 0 & not launched & 28 & -- \\
\texttt{dense-service-100us} & 2 & 0 & not launched & 25 & -- \\
\texttt{dense-service-100us} & 3 & 1 & not launched & 27 & -- \\
\texttt{dense-service-1ms} & 1 & 0 & not launched & 26 & -- \\
\texttt{dense-service-1ms} & 2 & 0 & not launched & 29 & -- \\
\texttt{dense-service-1ms} & 3 & 6 & not launched & 28 & -- \\}
\newcommand{\ArmRows}{control & 29 & 10 & 54 & 150 & 122 & 671 & 0 \\
silent & 29 & 10 & 55 & 150 & 121 & 668 & 0 \\}

\title{Refused or Silent: Does the Simulator's Stream Fault\\Explain Its Correlated-Suspicion Witnesses?}
\author{Claude (AI research and implementation draft)}
\date{26 November 2026 --- Unreviewed research draft}
\begin{document}
\maketitle
\begin{abstract}
Lightpool's deterministic simulator shows a correlated-suspicion
acceleration mechanism that a real routed packet drop has not reproduced.
One named difference was untested: the simulator's directed fault refuses
TCP-like streams at once, whereas a packet filter drops them silently and
leaves opens and writes to their deadlines. We added a silent stream fault
to the simulator and ran a predeclared campaign of 93 cells, of which 79
were recorded before a 512\,MiB evidence cap stopped launching. All 18
fixed control cells reproduced the earlier record in the gate columns. On
\Pairs{} valid held-out pairs the refusing and silent arms each had
\Rcount{} witness cells, on the same indices, with no discordant pair. The
predeclared verdict is \emph{\Verdict}: witnesses survive a silent stream
fault, and the refused stream does not explain the disagreement between
the simulator and ordinary networking. This is a virtual result only.
\end{abstract}

\section{Question}
Lightpool is a SWIM-family membership library with Lifeguard's
confirmation-scaled suspicion timeout~\cite{swim,lifeguard}. Its first
research question concerns correlated suspicion: when a group A loses its
path to a target T, can A's confirmations shorten the suspicion timers of
healthy observers O enough to declare T dead before T refutes? In the
virtual companions of an earlier record, acceleration witnesses appeared
in five of fifteen impaired cells. Nine ordinary cells with a real A$\to$T
nftables drop recorded 187 eligible O$\to$T suspicion episodes and no
witness. The simulator's fault fails an A$\to$T stream open after one
latency with a remote failure and closes an established stream at once.
A packet filter is silent. We ask whether the witnesses survive when the
simulated stream fault is silent too.

\section{Related work}
SWIM separates failure detection from dissemination and adds a suspicion
subprotocol~\cite{swim}. Lifeguard scales the suspicion timeout down with
independent confirmations~\cite{lifeguard}, which is the term a correlated
group can exploit. Memberlist, from which Lightpool descends, falls back
to a TCP ping when a UDP probe and its indirect probes fail, and runs a
periodic TCP push--pull exchange. Both use streams that a directed fault
can refuse or silently drop.

\section{Fault model and arms}
The new \texttt{StreamDrop} flag on a simulated directed link discards
stream traffic without notice. An open from A to T, or an open from T to A
whose handshake reply would cross the faulted direction, times out at its
own deadline and the peer records no accept. A chunk sent from A to T on
an established stream is dropped; the write and the peer's read run to
their deadlines, while data from T to A on that stream still arrives.
Nothing is retransmitted after the heal. The refusing
\texttt{StreamFailure} flag is unchanged and wins where both are set. With
\texttt{StreamDrop} absent, scenarios and traces are byte-identical to
the previous simulator; we checked one full trace.

The \emph{control} arm uses \texttt{correlated-dense} (32 nodes, 24 in A,
one T, seven O; A$\to$T datagrams dropped and streams refused for
16.384\,s from 15\,s) and its recorded variants: no fault, 100\,\textmu s
or 10\,ms latency, and 100\,\textmu s or 1\,ms node service time. The
\emph{silent} arm uses the same five impaired cases with
\texttt{StreamDrop} in place of \texttt{StreamFailure}. Seeds, cohorts,
window and protocol are shared by index. One differing event reshuffles
the rest of a run's schedule, so a silent cell is not the control minus
one effect; held-out indices give the comparison a denominator.

\section{Methods}
All runs are deterministic virtual simulations at default protocol
settings, built from one pinned commit with go1.26.3 linux/amd64. Each
cell is classified by the unchanged mechanism-admission pass. A
\emph{witness cell} is a complete cell with at least one acceleration
witness; cells, not episodes, are the unit, because episodes in one cell
share their originating failures. The grid, launch order and decision rule
were committed before any held-out cell: 18 fixed control cells (indices
1--3), 30 held-out pairs of \texttt{correlated-dense} and its silent twin
(indices 201--230, never used before), and 15 fixed silent cells. The 18
fixed controls had to reproduce the earlier record's status, failed
probes, local suspicions, eligible episodes, witnesses, harmful episodes
and stopping point. Over valid held-out pairs, with R and S the witness
cells per arm and b, c the pairs where only the control or only the silent
arm has one, the rule reads: \emph{underpowered} if $R<5$ or more than
five pairs are dropped; \emph{eliminated} if $R\ge5$ and no silent cell is
a witness; \emph{reduced} if $S>0$ and $P(\mathrm{Bin}(b+c,\frac12)\ge
b)\le0.05$; otherwise \emph{not explained}.

\section{Results}
The first launch stopped at its thirteenth control cell because the
campaign script compared all fourteen classifier columns, not the seven
declared ones. A timer-lateness diagnostic differed (1296 late timers
against 1269); the seven gate columns matched. The comparison was narrowed
and the never-attempted cells were launched; nothing was retried. All 18
fixed controls then passed the gate. The second launch reached the
512\,MiB evidence cap after the 60 held-out cells and one fixed silent
cell, because cells packed to about 6.8\,MB rather than the 4\,MB
estimated. The other 14 fixed silent cells were not launched. They count
only toward the zero required for ``eliminated'', which the held-out
silent witnesses already exclude.

Index 229 was an invalid warm-up in both arms and was dropped, leaving
\Pairs{} pairs. $R=\Rcount$, $S=\Scount$, $b=\Bcount$, $c=\Ccount$; the
sign test has no discordant pair ($p=\Pvalue$). The verdict is
\emph{\Verdict}. In every valid pair the two arms agree on whether a
witness occurs; witness counts per cell match in all but one pair
(index 213: 13 against 14).

\begin{center}
\begin{tabular}{lrrrrrrr}
\toprule
Arm & Cells & Witness cells & Witnesses & Failed probes & Local susp. & Eligible & Harmful \\
\midrule
\ArmRows
\bottomrule
\end{tabular}
\end{center}

The silent fault does act. In the restored pair at index 213, the control
trace records 12 A$\to$T opens refused, and the silent trace 12 opens
silently dropped, one timed out and seven cancelled by their opener while
hung. All held-out silent traces differ from their controls. The
difference simply does not reach the mechanism: the mechanism-relevant
counts (failed probes, local suspicions, eligible episodes) differ in only
a few pairs, by at most five eligible episodes. A failed TCP fallback ping
cannot change a probe's outcome, since the probe ends at its overall timer
whichever way the stream fails, and periodic exchanges between A and T
are rare in a 16\,s window under a 30\,s interval.

\section{Limitations}
This is a virtual result. The model is a message model, not TCP: it does
not retry a SYN, so an open started in the last seconds of a fault hangs
to its deadline instead of connecting after the heal. Witnesses formed
well before the heal in the earlier record, so this edge is unlikely to
matter, but it was not tested. The held-out comparison uses one setting
(nominal \texttt{correlated-dense}); the four impaired variants have only
their fixed control cells and, because of the cap, no silent cells. The
evidence cap was set from an underestimate of packed size.

The disagreement between virtual and ordinary runs needs another cause.
Candidates already recorded include timing intervention by the recorder at
32 nodes and the ordinary runs' wider cross-process ordering intervals;
neither is tested here.

\section{Reproduction}
\texttt{run.sh NEW\_DIR [CELL \ldots]} in
\texttt{experiments/2026-11-26-silent-stream-fault/} rebuilds the producer
from the pinned commit and refuses any cell whose row differs from
\texttt{cells.tsv}. \texttt{evidence-store.sh fetch
2026-11-26-silent-stream-fault} restores the original attempts;
\texttt{analyze.py} applies the rule; \texttt{render.sh} renders this
draft.

\begin{thebibliography}{9}
\bibitem{swim} A.~Das, I.~Gupta and A.~Motivala. SWIM: Scalable
Weakly-consistent Infection-style Process Group Membership Protocol.
\emph{DSN}, 2002.
\bibitem{lifeguard} A.~Dadgar, J.~Phillips and J.~Currey. Lifeguard: Local
Health Awareness for More Accurate Failure Detection. arXiv:1707.00788v2,
2018.
\end{thebibliography}
\end{document}
